\chapter{Against Love and Salvation}
At the end of Madoka Magica, Madoka ascends to a higher plane of being, sacrificing not only her life but her entire existence to save the other magical girls from becoming witches. As becomes clear from both Homura's explanation to Kyubey in the final episode of the series and comments by Nagisa and Sayaka in Rebellion, the magical girls so rescued continue existing in some form outside the universe, with Madoka. Whatever form they are in, we know they are in some sense aware and able to make decisions, and it appears are simultaneously magical girls and witches (which, of course, they always were).

That this is presented, within the series and initially within Rebellion, as a positive development and more-or-less happy ending is perhaps puzzling. Arguably the entire point of Sayaka's character arc was coming to realize that it was a mistake to try to sacrifice herself to save another, while Homura's attempts to save Madoka were similarly depicted as making things continually worse for them both. It is not particularly surprising, then, that Rebellion calls that salvation into question.

The first segment of the movie, in which the magical girls are happy and get along, and the opponents they face are challenging but conquerable, serves as a parody of both worlds that can be understood as "Madoka's world." As a new enemy representing human misery, the Nightmares are a twisted reflection of the Wraiths. Like the Wraiths, the Nightmares are oddly similar to one another, but where the Wraiths are fairly creepy, attenuated humanoid giants, the Nightmares wear bear suits and fire stuffed animals from their arms. Defeating a Wraith earned many small rewards for cleaning the magical girls' Soul Gems, making magical girl teams viable, unlike in the prior, witch-infested timelines. But in Homura's fantasy world, defeating a Nightmare creates a diffuse glow that purifies the Soul Gems, making magical girl teams actively preferable to intercept more of that light. In addition, where the first we see of the Wraith world created by Madoka is the death of Sayaka and mourning of her teammates, Homura's dream world preserves both Sayaka's life and her wish, by pairing (or at least heavily implying that it pairs) her with Kyoko to allow Hitomi and Kyousuke to be together. Homura's fantasy world is, simply, happier than the one Madoka created!

It is also sillier, and not just because of the bear suits. The juxtaposition of the mundane and the eerie is the province of surrealist art, and it is here that the Wraiths, and by extension the city they haunt, largely fall. Men with missing or obscured faces are, for instance, a favored subject of Magritte \cite{ref116}. Likewise, the witches, though more playful, are juxtaposed with extreme violence, both by the witches against humans and magical girls, and by magical girls against the witches. This combination of playful, often childish, imagery and violence forms a sort of brutalist surrealism. 

There is, however, no violence where the Nightmares are concerned. They destroy property, seemingly, but there is no trace of damage when the magical girls are done, and against them the magical girls are shown deploying traps and bindings or firing weapons to drive the Nightmare into a trap, but never attacking the Nightmare directly. The actual defeat of the Nightmare seems to involve actions that are at once highly ritualized, yet seemingly arbitrary—a banquet catered by the magical girls in the cold open, and a nursery rhyme-like chant or game about food against Hitomi's Nightmare. 

Oddly harmless violence, food, rhythmic verse, games, arbitrary rituals—these are all common features of nonsense literature \cite{ref117}. At the core of nonsense is an interest in alternative logics, in circumstances (such as games, meals, etiquette) where ultimately arbitrary, yet internally consistent, rules guide behavior; like a dream, nonsense substitutes one set of arbitrary rules for another, and lets the consequences play out logically \cite{ref118}. And yet within this nonsense, all five magical girls are alive and happy and thriving; it seems a world of nonsense is better than the world of Wraiths Madoka created. 

Madoka's Pure Land, her heaven, is also depicted as inferior. As I argued in Interlude 2, Madoka's "afterlife" is inferior to even Homura's dream world because it is a deathless world that contains no decay, no suffering, no putrefaction; both Sayaka and Nagisa chose to reify themselves alongside Madoka because they sought something that only existed as a consequence of decay and death, namely Sayaka's relationship with Kyoko and, for Nagisa, cheese. 

Homura's dream world is also more directly a parody of Madoka's "heaven," in the sense that Homura snatched magical girls (as well as a handful of ordinary humans) into her world without their consent and now keeps them there, trapped and cut off from the universe, but artificially happy. She has "saved" them because she has grown to care about them by extension, as the people Madoka loved—and at least in the case of Kyoko and Mami, whom she ultimately trusts to kill Homulilly, come to respect and possibly even like them, as well.

To want to save someone is necessarily to want power over that someone. By becoming a knight protector, Sayaka made herself a judge (and in the case of those two misogynists on the train, likely executioner as well). By wishing to be Madoka's protector, Homura ultimately put herself in a position to repeatedly try to take the choice of becoming a magical girl away from Madoka. And by wishing to save all magical girls from their destiny of becoming what they fight, Madoka set herself up as a goddess. 

To be a savior (as always, as opposed to helping, which involves the consent of the one helped and places the helper in a temporarily subordinate, rather than dominant, position), in other words, necessarily entails being a little bit of a tyrant. Since the savior is acting without the consent of the saved, they are very likely to get it wrong, as Madoka does with Homura. Look at the opening credits: Homura is depicted as a grey, troll-like figure kneeling while the magical girls dance. She is not capable of joining their happiness; the closest she is able to come is as the weak and shy "pigtails" version of her character during the first segment of the movie, and even then she is able to sense that something is deeply wrong. Once her hair is again loose, she is never genuinely happy again for the rest of the movie, for the simple reason that her untold ages of suffering, and the fact that she and she alone remembers them, have warped her emotionally to the point that she very possibly cannot be saved. 

Instead, she acts in parody of Madoka, snatching people up and placing them in her labyrinth. But is it really any different from what Madoka did? Is Madoka's sacrifice an act of selfless love while Homura's is selfish? And which is the greater sacrifice—your existence or your soul? Is it worse to never have existed, or to become the enemy of all you once held dear?

The answer, of course, is that it's a silly question. All value is relative, so it is entirely a matter of perspective which is worse; very likely, each of the two girls feels their own sacrifice is the greatest they could make, since Madoka cares deeply about her connections to others, while Homura is more focused on her cause. 

But, seeing in Homura's actions a twisted reflection of Madoka's, we see Madoka's in a new way as well. Can an act truly be considered selfless if it gets you everything you ever wanted? Madoka gets to be with, in her own words, "everyone"; all her loved ones are safe; she gets to defeat all the witches; she gets to become a magical girl; she gets to matter, quite possibly more than anyone else who ever lived. By contrast, Homura's choice to become a "demon" devoted to keeping Madoka in the world costs her the only thing she values, the chance to be together with Madoka in the end; now they must eventually be enemies. Isn't it therefore Homura who is selfless?

Of course not, because selfless love is an oxymoron. That is the point in depicting Homura's possessiveness, and through it revealing Madoka's selfishness. To love someone is to want to protect that person, possibly from themselves. It is to want to spend time with that person. It is to want that person to want you. Expressed in a healthy way and reciprocated equally, of course, love can be a wonderful thing; romantic or otherwise, it is the ultimate bond between two people. But like any bond, it can also be used to entrap, to control, to assert dominance—hence the existence of abusive relationships. 

Throughout the series, we saw magical girls torn between acknowledging what they genuinely wanted and what they believed they should want. Mami tortured herself for wishing to live, rather than wishing to save her parents. Sayaka and Kyoko wished for others' benefit, rather than wishing for those others to appreciate the help, and suffered tremendously as a result. This is why Kyubey targets girls, because from the moment they are given their first doll they are indoctrinated to take care of others, socialized to think of themselves as caretakers, responsible for the wellbeing of others \cite{ref119}. Society has done Kyubey's work for him, creating girls who will wish for what social pressure tells them they want instead of truly wishing for what they desire. (Not that it matters in the end, of course; the wish alone damns the magical girl to become a witch or die, though a poorly chosen wish makes the hope-despair cycle faster.)

So Homura sees no way to wish for what she truly desires, to be with Madoka. She wishes instead to take care of Madoka, first in the series at the end of the "original" timeline shown in the first part of Episode 10, and then in Rebellion when she becomes a "demon." In both cases, she ultimately sees no hope but to become "evil." Rebellion thus closes the largest cycle in a series full of cycles: the evolution of Homura Akemi from a dark, seemingly villainous character who disrupts the status quo to a dark, seemingly villainous character who maintains the status quo. More than ever, she is now Mami's dark mirror.
